\chapter{Conclusion}\label{conclusion}

\section{Summary}\label{summary}

This thesis investigated automated verification of textual UI requirements
against ordered screenshot flows. Its output contract combines a fulfillment
label with identifiable UI evidence and preserves uncertainty when the recorded
flow is insufficient.

To study this task, the thesis introduced an explicit four-label verification
contract, separated UI evaluability from fulfillment, implemented a modular
evidence-first pipeline, and constructed a reviewed 258-item benchmark over 13
Mind2Web-derived flows. The evaluation compared models, claim policies, and
screenshot-selection strategies while reporting label quality, unsafe positive
predictions, abstention, traceability, stability, cost, and qualitative error
patterns. A region-grounding extension adds spatially precise evidence without
being treated as a cause of improved label accuracy.

\section{Answers to the Research Questions}\label{answers-to-the-research-questions}

For \textbf{RQ1}, multimodal models assess requirement fulfillment from ordered
screenshot flows with useful but incomplete accuracy under the fixed
four-label operationalization. Gemini 3.1 Flash-Lite reaches 79.5\% accuracy and
0.511 macro-F1 in the matched raw/all configuration, compared with 73.3\% and
0.413 for Gemini 2.5 Flash-Lite. The hosted Qwen baseline shows that equal
headline accuracy can conceal substantially different false-fulfillment and
minority-class behavior. Model choice therefore matters, and accuracy alone
does not characterize fulfillment-label performance.

In the order-unavailable robustness condition, 10.9\% of labels change and aggregate accuracy decreases
by about 4.7 percentage points. Ordered input is therefore operationally
relevant to the implemented verifier, although the limited subgroup evidence
does not support a universal causal claim for every temporal requirement
pattern.

For \textbf{RQ2}, the effects of decomposition and screenshot selection are
conditional. Gated automatic decomposition does not materially improve
all-screenshot accuracy and increases false fulfillment in that setting, but it
improves macro-F1 when evidence is restricted to top-4. Shared top-4 selection
reduces raw-requirement accuracy and chronological evidence-trace alignment relative to the complete flow
while providing negligible cost savings in the current batching
implementation. Explicit evidence remains valuable for traceability even where
the staged configuration does not improve labels. Region-level evidence is
evaluated as an additional traceability output rather than an accuracy factor.

For \textbf{RQ3}, the main failures are systematic. Models over-generalize from a
visible entry point to an unobserved outcome, infer hidden or external behavior
from interface proxies, overstate universal and comparative claims, and miss
decisive late states when retrieval favors lexical overlap.

Native abstentions often protect against unsupported closed-world decisions:
replacing all abstentions with negative labels substantially reduces accuracy.
This does not prove that every abstention is calibrated, but it confirms that
insufficient evidence and visible contradiction must remain distinct.

The completed
screenshot-aware audit assigns 142 of 392 label errors (36.2\%) to unsafe
over-fulfillment, 138 (35.2\%) to excessive abstention, and 57 (14.5\%) to
evidence-selection misses. Multi-screen composition, late result states, hidden
or external behavior, and cross-step persistence are the most frequent
overlapping requirement patterns. The results answer RQ3 descriptively for the
frozen 13-flow benchmark; repeated conditions and the LLM-assisted coding with
targeted primary-author boundary review do not support population-level
prevalence or independent-human reliability claims.

\section{Contributions}\label{contributions}

The work defines evidence-bounded UI requirement verification over ordered
observations and makes its visible/hidden boundary and label policy explicit. A
modular implementation covers flow ingestion, requirement understanding,
screenshot selection, multimodal claim verification, deterministic
aggregation, evidence localization, annotation, and experiment packaging.

The accompanying benchmark records requirement labels, UI evaluability,
claims, evidence steps, rationales, uncertainty reasons, and contrastive items.
Its reference annotations were produced through prediction-independent
primary-author review. The controlled experiments reveal an
interaction between decomposition and retrieval and show that smaller
screenshot sets can omit decisive evidence without meaningful cost savings.
The evaluation framework keeps screenshot retrieval, region grounding, UI
evaluability, abstention, runtime, cost, and stability as distinct outcomes.

\section{Future Work}\label{future-work}

The next methodological step is an error-driven refinement loop for the label
policy. After each evaluation, the primary error categories can be inspected
for recurring mechanisms and translated into general evidence-sufficiency
rules. For example, the rubric could distinguish visible states from downstream
effects, bounded UI sets from open-world completeness, and strong visible
proxies from hidden guarantees. The revised guide, prompt, and deterministic
gates would then be evaluated on held-out flows, and the remaining errors would
start the next iteration. This process could reduce model-dependent assumptions
without treating the observed categories as universal causes or assuming that
each iteration must improve accuracy.

This work requires a larger, requirements-first benchmark. Independent
requirements should be fixed before collecting corresponding UI drafts,
implementations, and execution flows, ideally with several variants containing
natural or deliberately introduced deviations. Such data would separate
requirement ambiguity, design conformance, implementation defects, and
incomplete execution evidence. The rubric should be developed separately from
the evaluation applications and tested on held-out contexts. More diverse
interfaces, requirements from real projects, additional flow clusters, and
several reviewers would also strengthen estimates for rare classes and reveal
whether refined rules transfer beyond the current benchmark.

For short flows, future systems should retain direct complete-flow verification
as the baseline because neither mandatory claim decomposition nor shared
lexical top-k selection improved the strongest configuration here. Adaptive evidence
acquisition becomes relevant when flows exceed practical context limits. It
should pair actions with result states, prioritize late states, and request more
screenshots when evidence is insufficient. Finally, region grounding should
improve proposal coverage and support multi-region or transition evidence. Its
success should be measured through relevance, sufficiency, reviewer effort, and
auditability rather than assumed to improve the verification label.
